\documentclass[11pt]{article}
\usepackage[margin=1in]{geometry}
\usepackage[hidelinks]{hyperref}
\setlength{\parindent}{0pt}
\setlength{\parskip}{0.7em}
\pagestyle{empty}

\begin{document}

\noindent Tomas Rojas\\
Sede de Occidente, Universidad de Costa Rica, San Ram\'on, Costa Rica\\
Centro de Investigaci\'on en Ciencia e Ingenier\'ia de Materiales (CICIMA),\\
Universidad de Costa Rica, San Jos\'e, Costa Rica\\
\texttt{tomas.rojas\_s@ucr.ac.cr}

\vspace{1em}
\noindent 17 July 2026

\vspace{1em}
\noindent To the Editors,\\
\textit{Physica E: Low-dimensional Systems and Nanostructures}

\vspace{1em}
\noindent Dear Editors,

We are pleased to submit our manuscript, ``Thermoelectric transport and spin filtering in
hexagonal CrN nanoribbons: a tight-binding Landauer study,'' for consideration as a research
article in \textit{Physica E}.

Chromium nitride is an established thermoelectric material in its bulk and thin-film forms, and
its two-dimensional hexagonal allotrope has attracted sustained interest as a half-metallic
magnet---yet the thermoelectric response of CrN in the nanoribbon geometry has never been
reported. Our manuscript closes that gap with a deliberately lightweight approach: a
Slater--Koster tight-binding model parametrized entirely from published first-principles data
(bands and phonons), combined with phase-coherent Landauer transport.

The paper makes three contributions. First, it provides the first spin-resolved thermoelectric
characterization of CrN nanoribbons: honest figures of merit of $ZT\simeq0.01$--$0.15$ at room
temperature, with the global optimum pinned to the minority conduction-band edge, and with one
geometry (the narrow armchair ribbon) uniquely combining its best $ZT$ with $100\%$ spin
polarization. Second, it establishes that these ribbons are outstanding, gate-robust spin
filters---the minority channel is gapped over more than $4$~eV---and shows that the same gap
powers a thermally driven spin valve: for collinear antiparallel junctions the coherent
transmission, and hence the thermocurrent, vanishes identically across the window, and a
noncollinear domain-wall calculation quantifies the required sharpness of the magnetization
reversal, singling out geometrically constrained or spacer-separated junctions. The magnetic edges and
their exchange-coupling hierarchy are reproduced self-consistently. Third, and of methodological
interest to the low-dimensional transport community, it documents quantitatively how the
\emph{minimal} orbital manifold suggested by the states at the Fermi level fabricates a sharp
transmission edge and overestimates $ZT$ by almost an order of magnitude---an artifact that
parameter-sensitivity analysis cannot detect, and a concrete caution for minimal-basis
nanoribbon thermoelectric studies in general.

We believe the manuscript fits the scope of \textit{Physica E} well: it concerns electronic,
magnetic and thermoelectric properties of a nanoribbon system, uses transport methodology
standard in the low-dimensional community, and is complementary---in property and in method---to
the existing first-principles/DMRG work on CrN nanoribbon magnetism.

This manuscript is original, has not been published previously, and is not under consideration
elsewhere. All authors have approved the submission and declare no competing interests. The
code, digitized reference data and scripts that generate every figure are provided in an
accompanying open repository to support full reproducibility.

Thank you for considering our manuscript.

\vspace{1em}
\noindent Sincerely,\\
Tomas Rojas\\
\textit{on behalf of the authors (T. Rojas and R. Thorat)}\\
Corresponding author

\end{document}
